\documentclass[11pt]{article}
\usepackage[a4paper,margin=27mm]{geometry}
\usepackage[T1]{fontenc}
\usepackage[utf8]{inputenc}
\usepackage{lmodern,amsmath,amssymb,amsthm,mathtools,microtype,booktabs,array}
\usepackage{hyperref}
\hypersetup{colorlinks=true,linkcolor=blue,citecolor=blue,urlcolor=blue}
\newtheorem{theorem}{Theorem}
\newtheorem{lemma}[theorem]{Lemma}
\newtheorem{proposition}[theorem]{Proposition}
\newtheorem{corollary}[theorem]{Corollary}
\theoremstyle{remark}\newtheorem{remark}[theorem]{Remark}
\DeclareMathOperator{\diag}{diag}
\DeclareMathOperator{\Hess}{Hess}
\DeclareMathOperator{\rank}{rank}
\newcommand{\R}{\mathbb R}
\newcommand{\Q}{\mathbb Q}
\newcommand{\one}{\mathbf 1}
\newcommand{\cI}{\mathcal I}
\title{Six-spin Lorentzian obstructions for clique blowups\\of independence-stable graphs}
\author{Lluis Eriksson\\\small Independent researcher, Sweden}
\date{17 October 2026\\\small Working-paper manuscript}
\begin{document}
\maketitle
\begin{abstract}
Let $F$ be a connected simple graph with at least one edge, and suppose
that deleting any edge leaves its independence number unchanged.
For every integer $s\ge2$, we construct a strictly positive rational
antiferromagnetic matrix of rank two on six spins for which the
partition polynomial of the clique blowup $F[K_s]$ is not Lorentzian.
The number of spins is independent of the order and degrees of $F$
and of the blowup size. The matrix itself may depend on $F$ and $s$.
The proof combines the classical equality case of Turan's theorem
with a conditioned Hessian identity. Edge-deletion independence
stability forces all minimum-edge occupations onto maximum
independent sets. The first four active profiles then form a
nonsingular block of inertia $(2,2,0)$; a two-stage limiting argument
realizes this block using six spins. The theorem includes clique
blowups of all connected biregular bipartite graphs with at least
one edge other than $K_2$, and of connected graphs with at least
one edge and a unique maximum independent set.
Exact checks cover 720 blowups of 360 bases of orders three through
seven, and four explicit positive six-spin witnesses.
The connected-source classification conjecture remains open.
\end{abstract}

\section{Statement and relation to earlier work}

For a finite simple graph $G=(V,E)$ and an entrywise nonnegative
symmetric matrix $K$ indexed by $q$ spins, define
\begin{equation}
 Z_G(\mathbf x;K)=
 \sum_{\sigma:V\to\{0,\ldots,q-1\}}
 \prod_{uv\in E}K_{\sigma(u),\sigma(v)}
 \prod_{v\in V}x_{\sigma(v)}.                    \label{eq:partition}
\end{equation}
We use the spectral convention that $K$ is \emph{antiferromagnetic}
when it has at most one positive eigenvalue. Lee--Oh--Seo prove that
clique partition polynomials are Lorentzian for every such target
and propose a converse \cite{LOS}. We use the connected-source
formulation discussed by Lee--Seo \cite{LS}: must a connected graph
whose partition polynomial is Lorentzian for every antiferromagnetic
target be a clique? Connectedness is necessary because partition
polynomials factor over components and Lorentzian polynomials are
closed under products \cite{BH}.

For a homogeneous polynomial of degree $n\ge2$ with strictly positive
coefficients on the full degree-$n$ simplex, Lorentzianity is
equivalent to every partial derivative of total order $n-2$ having
a Hessian with at most one positive eigenvalue
\cite[Theorem 2.25]{BH}. The partition polynomials for positive
targets have this full support. Inertia triples list positive,
negative and zero eigenvalue counts, in that order.

Write $\alpha(F)$ for the independence number. A graph is
\emph{$\alpha^-$-stable} if $\alpha(F-e)=\alpha(F)$ for every edge $e$.
This terminology and its characterization by maximum independent
sets are classical \cite[Theorem 2.1]{LM}. The graph $F[K_s]$ replaces
each vertex of $F$ by a clique of order $s$, and each edge by all
edges between the corresponding cliques.

\newpage
\begin{theorem}\label{thm:main}
Let $F$ be a connected $\alpha^-$-stable simple graph with at least
one edge. For every integer $s\ge2$, there are distinct rational
numbers $a_0,\ldots,a_5$ in $(-1,1)$, with $a_1=0$, such that
\[
 K=\one\one^\top-aa^\top
\]
is strictly positive entrywise and
$Z_{F[K_s]}(\mathbf x;K)$ is not Lorentzian.
The matrix $K$ has inertia $(1,1,4)$.
\end{theorem}

The six-spin bound is uniform. The rational parameters are not
asserted to be uniform, and the theorem does not identify the
smallest possible number of spins.

The preceding work \cite{Eriksson} establishes a positive rank-two
obstruction for connected noncliques with no true twins among their
maximum-degree vertices, together with several other source classes.
Every vertex of $F[K_s]$ has true twins when $s\ge2$, so that
maximum-degree theorem does not apply to these blowups. Some special
families overlap earlier results, including particular joins of
cliques. The present statement supplies one criterion for all
$\alpha^-$-stable bases and a spin bound independent of graph size.
It is not a claim to solve the full conjecture.

For a concrete distinction, $C_6[K_2]$ is covered here. It is
$5$-regular with true twins, and adjacent pairs have either two
or four common neighbours, so it is not edge-regular. The earlier
neighbourhood-gap quantities satisfy $b=c=2$, so that strict-gap
criterion also does not apply. In contrast, $P_3[K_s]$ overlaps
the earlier join-of-cliques family. The abstract block form is
already familiar from \cite{Eriksson}; the contribution here is
its uniform occurrence at the \emph{first} active profiles, together
with suppression of every later profile by six-spin compression.

\begin{center}
\begin{tabular}{@{}>{\raggedright\arraybackslash}p{.28\textwidth}>{\raggedright\arraybackslash}p{.34\textwidth}>{\raggedright\arraybackslash}p{.28\textwidth}@{}}
\toprule
Ingredient & Role here & Attribution\\
\midrule
$\alpha^-$-stability & Controls maximum independent sets &
Classical \cite{LM}\\
Turan equality & Classifies the leading occupations &
Classical \cite{Turan}\\
Conditioned profiles & Expresses the Hessian &
Earlier framework \cite{Eriksson}\\
Initial four-profile block & Produces two positive directions &
Present argument\\
Six-spin compression & Realizes those directions in a positive target &
Present argument\\
\bottomrule
\end{tabular}
\end{center}

\section{Conditioned Hessians and their leading term}

We recall the identity needed from \cite{Eriksson}, including its
derivation to make the argument self-contained. Choose target nodes
$a_0=\ell$, $a_1=0$, and $a_i\in(-1,1)$, where $\ell\ne0$, and put
$\theta=\ell^2$. Let $n=|V(G)|$ and $0\le h\le n-2$.

For an $h$-vertex subset $S$, define the profile of an outside vertex by
\[
 r_S(v)=|N_G(v)\cap S|.
\]
Let $\nu^S_r$ count outside vertices with profile $r$, and let
$C^S_{rt}$ count \emph{ordered} adjacent pairs of distinct outside
vertices with respective profiles $r,t$. Thus $C^S$ is symmetric.
Pad $\nu^S$ with a zero in coordinate $h+1$. Let
$D e_r=e_r-e_{r+1}$, for $0\le r\le h$. Here $C^S$ is
$(h+1)\times(h+1)$ and $D$ is $(h+2)\times(h+1)$. Set
\begin{equation}
 Q_h(\theta)=\sum_{|S|=h}(1-\theta)^{e(G[S])}
 \left[
 \theta\bigl(\nu^S(\nu^S)^\top-\diag\nu^S\bigr)
 -D C^S D^\top
 \right].                                      \label{eq:Q}
\end{equation}
The matrix has profile coordinates $0,\ldots,h+1$.

\begin{lemma}\label{lem:identity}
For $U_{ir}=(1-\ell a_i)^r$,
\begin{equation}
 \Hess\!\left(\partial_0^h\partial_1^{n-h-2}
       Z_G(\mathbf x;K)\right)
 =\frac{h!(n-h-2)!}{\theta}\,
       UQ_h(\theta)U^\top.                       \label{eq:hessian}
\end{equation}
This identity holds for any number of target states.
\end{lemma}
\begin{proof}
Choose the $h$ vertices differentiated in spin zero; they form $S$.
All but two of the other vertices are differentiated in spin one.
Because $a_1=0$, all interactions involving those latter vertices
have weight one. The internal edges of $S$ contribute
$(1-\theta)^{e(G[S])}$.
If the two remaining vertices have profiles $r,t$ and spins $i,j$,
their interactions with $S$ contribute
$(1-\ell a_i)^r(1-\ell a_j)^t$.
Their mutual edge, when present, adds the factor $1-a_i a_j$.
The matrix $\nu\nu^\top-\diag\nu$ counts all ordered distinct pairs,
and $C$ counts the ordered edges.
Finally
\[
 (UD)_{ir}=(1-\ell a_i)^r-(1-\ell a_i)^{r+1}
           =\ell a_i(1-\ell a_i)^r.
\]
This yields \eqref{eq:hessian}, including the differentiation
factorials.
\end{proof}

A Lorentzian polynomial has at most one positive Hessian eigenvalue
in every quadratic partial derivative \cite{BH}. Thus two positive
eigenvalues in any principal restriction of \eqref{eq:hessian}
provide an obstruction.

Let
\[
 e_0=\min_{|S|=h}e(G[S]).
\]
The polynomial identity \eqref{eq:Q} gives
\begin{equation}
 \lim_{\theta\uparrow1}
 \frac{Q_h(\theta)}{(1-\theta)^{e_0}}=L,\qquad
 L=\sum_{\substack{|S|=h\\ e(G[S])=e_0}}
 \left[\nu^S(\nu^S)^\top-\diag\nu^S-D C^S D^\top\right].
                                                        \label{eq:leading}
\end{equation}
All limits are in fixed finite-dimensional matrix spaces. We call
$L$ the leading profile kernel. A profile index is \emph{active}
if its row in $L$ contains a nonzero entry. It is a rescaled limit; the
unscaled derivative at $\ell=-1$ may vanish.

\section{Minimum-edge occupations}

We first record the classical stability characterization and the
part used in the proof.

\begin{lemma}[Classical; see \cite{LM}]\label{lem:stability}
The following are equivalent:
\begin{enumerate}
\item $\alpha(F-e)=\alpha(F)$ for every edge $e$;
\item every maximum independent set $I$ is $2$-dominating:
$|N_F(u)\cap I|\ge2$ for every $u\notin I$.
\end{enumerate}
Under these conditions, two maximum independent sets cannot differ
by replacing a single vertex.
\end{lemma}
\begin{proof}
Maximality excludes zero neighbours in $I$. If $u$ has exactly one
neighbour $v\in I$, deleting $uv$ makes $I\cup\{u\}$ independent.
Conversely, an independent $(\alpha(F)+1)$-set created by deleting
$uv$ must contain both endpoints; removing $u$ gives a maximum
independent set in which $u$ has just the neighbour $v$.
For the last assertion, a replacement vertex is nonadjacent to all
the retained vertices, and can have at most one neighbour in the
original maximum independent set.
\end{proof}

Let $\alpha=\alpha(F)$, $G=F[K_s]$, and
\begin{equation}
 h=\alpha s-1,\qquad
 e_0=(\alpha-1)\binom{s}{2}+\binom{s-1}{2}.        \label{eq:h}
\end{equation}
The hypotheses of Theorem~\ref{thm:main} imply $\alpha\ge2$:
a connected graph with an edge and $\alpha=1$ is complete and is
not $\alpha^-$-stable. Also $\alpha<|V(F)|$, so $h\le |V(G)|-2$.

\begin{lemma}\label{lem:ground}
The minimum number of edges induced by an $h$-vertex subset of $G$
is $e_0$. Every minimizing subset is obtained from a maximum
independent set $I$ of $F$ by selecting all vertices of its classes
except for one omitted vertex in one class $v\in I$.
For each pair $(I,v)$ there are exactly $s$ such subsets.
\end{lemma}
\begin{proof}
The independence number of every induced subgraph of $G$ is at
most $\alpha$. Turan's theorem and its equality case
\cite{Turan}, applied to the complement of an $h$-vertex induced
subgraph, imply the bound \eqref{eq:h}. Equality requires a
disjoint union of $\alpha$ cliques: one of order $s-1$ and the
others of order $s$. The construction in the statement attains it.

It remains to exclude other equality occupations. Choose one
occupied base class from each of these $\alpha$ clique components.
Their base vertices form a maximum independent set in $F$.
If a component occupied two or more base classes, replacing its
chosen representative by another would give two maximum
independent sets differing in one vertex. This contradicts
Lemma~\ref{lem:stability}. Therefore each component occupies one
class, and its size determines the asserted construction.
The omitted vertex has $s$ choices.
\end{proof}

This argument uses the equality case of a classical extremal
theorem. It does not assert a new extremal edge bound.

\section{The initial four-profile block}

Let $\cI$ be the set of maximum independent sets of $F$.
For $I\in\cI$ and $u\notin I$, put $d_I(u)=|N_F(u)\cap I|$, and define
\begin{equation}
 N_d=\bigl|\{(I,u): I\in\cI,\ u\notin I,\ d_I(u)=d\}\bigr|,
 \qquad d_0=\min\{d:N_d>0\}.                    \label{eq:Nd}
\end{equation}
By Lemma~\ref{lem:stability}, $d_0\ge2$.

For a minimizing occupation indexed by $(I,v)$, the single remaining
vertex of class $v$ has profile $s-1$. Each outside class $u\notin I$
has $s$ remaining vertices and profile
\begin{equation}
 r_u=s\,d_I(u)-\mathbf 1_{\{uv\in E(F)\}}\ge2s-1.
                                                        \label{eq:profiles}
\end{equation}
All other classes in $I$ have no remaining vertices.

\begin{proposition}\label{prop:block}
The first four active profile coordinates of $L$ are
\[
 A=(s-1,\ s,\ d_0s-1,\ d_0s).
\]
The corresponding principal block has the form
\begin{equation}
 L[A,A]=
 \begin{pmatrix}0_{2\times2}&B\\B^\top&H\end{pmatrix},
 \qquad
 B=s^2N_{d_0}
 \begin{pmatrix}0&\alpha\\d_0&-d_0\end{pmatrix}. \label{eq:block}
\end{equation}
It has inertia $(2,2,0)$, and
\[
 \det L[A,A]=s^8\alpha^2d_0^2N_{d_0}^4>0.
\]
\end{proposition}
\begin{proof}
There is exactly one remaining vertex of profile $s-1$, no
remaining vertex of profile $s$, and no earlier profile.
The internal ordered-edge contribution at $s-1$ is zero.
Consequently the leading $2\times2$ block on $s-1,s$ is zero.

We compute its couplings to the other coordinates. For a fixed
$(I,v)$ and an outside class $u$, write $d=d_I(u)$.
The term $\nu\nu^\top-\diag\nu$ contributes $s$ between the
remaining vertex of class $v$ and class $u$.
If $uv$ is an edge, the edge correction cancels this contribution
at $ds-1$ and restores it at $ds$ through the difference operator.
If $uv$ is not an edge, its profile is already $ds$.
The row at $s$ is the opposite difference of the two edge endpoints.
Summing over $v\in I$ and the $s$ choices of omitted vertex yields
\begin{align}
 L_{s-1,ds}&=s^2\alpha N_d,&
 L_{s-1,ds-1}&=0, \label{eq:row1}\\
 L_{s,ds-1}&=s^2dN_d,&
 L_{s,ds}&=-s^2dN_d. \label{eq:row2}
\end{align}
All other entries in these two rows vanish. Edges between outside
classes affect only the remaining block, and need not be counted.

Equation~\eqref{eq:profiles} and the one-step difference operator
show that no coordinate strictly between $s$ and $d_0s-1$ is active.
The four displayed coordinates are distinct because $s,d_0\ge2$;
their couplings show that they are active. This proves
\eqref{eq:block}. Its off-diagonal block has determinant
$-s^4\alpha d_0N_{d_0}^2\ne0$.
An invertible congruence removes $H$, leaving
$\left(\begin{smallmatrix}0&B\\B^\top&0\end{smallmatrix}\right)$,
which has two eigenvalues of each sign.
The determinant is $(\det B)^2$.
\end{proof}

\section{Compressing to six spins}

The fact that the block in Proposition~\ref{prop:block} uses the
\emph{first} active profiles is essential to the following argument.

\begin{lemma}\label{lem:compression}
Let $L$ be a finite real symmetric matrix indexed by distinct
nonnegative integers. Let $b_1<\cdots<b_r$ be its first $r$ active
indices, and suppose $L[B,B]$ is nonsingular. Choose distinct
positive numbers $z_1,\ldots,z_r$, and put
$U_\delta(i,j)=(\delta z_i)^j$ on the active columns.
For sufficiently small $\delta>0$, $U_\delta L U_\delta^\top$
has the same inertia as $L[B,B]$.
\end{lemma}
\begin{proof}
Let $V_\delta=U_\delta[:,B]$. The generalized Vandermonde matrix
$V_1=(z_i^{b_j})$ is invertible. Indeed, a nonzero linear
combination of $r$ distinct monomials has at most $r-1$ distinct
positive zeros, by Descartes' rule of signs.
Since $V_\delta=V_1\diag(\delta^{b_j})$, every entry of
$V_\delta^{-1}U_\delta[:,t]$ for $t>b_r$ is a constant times
$\delta^{t-b_j}$, and tends to zero. Hence
\[
 V_\delta^{-1}U_\delta L U_\delta^\top V_\delta^{-\top}
 \longrightarrow L[B,B].
\]
Inertia is locally constant at a nonsingular real symmetric matrix,
and invariant under invertible congruence.
\end{proof}

\begin{proof}[Proof of Theorem~\ref{thm:main}]
Apply Lemma~\ref{lem:compression} with $r=4$,
$z_i=i$, and the block in Proposition~\ref{prop:block}.
Choose a sufficiently small rational $\delta\in(0,1/4)$ so that
$U_\delta L U_\delta^\top$ has inertia $(2,2,0)$.

Fix this $\delta$, and choose the six target nodes
\begin{equation}
 a_0=\ell,\qquad a_1=0,\qquad
 a_{i+1}=-1+\delta i\quad(1\le i\le4).           \label{eq:nodes}
\end{equation}
As $\ell\downarrow-1$ from above,
$1-\ell a_{i+1}\to\delta i$.
Let $W_\ell$ consist of rows $2,3,4,5$ of $U$.
By \eqref{eq:leading},
\[
 \frac{W_\ell Q_h(\ell^2)W_\ell^\top}{(1-\ell^2)^{e_0}}
 \longrightarrow U_\delta L U_\delta^\top.
\]
It retains two positive eigenvalues for $\ell>-1$ sufficiently
close to $-1$. Choose such an $\ell$ rational, and sufficiently
close that it is different from all other nodes.

All six nodes now lie in $(-1,1)$, so every entry of $K$ is
strictly positive. The vectors $\one$ and $a$ are independent;
the representation
\[
 K=\begin{pmatrix}\one&a\end{pmatrix}
       \begin{pmatrix}1&0\\0&-1\end{pmatrix}
       \begin{pmatrix}\one&a\end{pmatrix}^{\!\top}
\]
gives inertia $(1,1,4)$.
The positive scalar in \eqref{eq:hessian} preserves the two
positive eigenvalues of the principal Hessian restriction.
Therefore the quadratic partial derivative is not Lorentzian,
and neither is $Z_G$.
\end{proof}

The limits are taken in a specified order: first choose $\delta$,
then choose $\ell$. No interchange or graph-uniform neighbourhood
is used. In particular, the argument does not assert that the
boundary target with $\ell=-1$ is itself an obstruction.

\section{Examples and a finite construction procedure}

\begin{corollary}\label{cor:biregular}
Theorem~\ref{thm:main} applies to every connected biregular
bipartite base with at least one edge other than $K_2$.
\end{corollary}
\begin{proof}
Let the bipartition sizes be $a\le b$, with respective degrees
$D\ge d>0$, so $aD=bd$. For an independent set $J$,
\[
 D|J\cap A|+d|J\cap B|\le |E(F)|=db.
\]
It follows that $\alpha(F)=b$. If $D>d$, equality in cardinality
forces $J=B$, so the maximum independent set is unique.
If $D=d$, equality means every edge has exactly one endpoint in
$J$, so the complement is also independent. Connectedness gives
just the two bipartition classes. Except for $K_2$, those two
classes cannot differ by a single-vertex replacement.
In either case a maximum independent set cannot have an outside
vertex with just one neighbour in it; apply Lemma~\ref{lem:stability}.
\end{proof}

\begin{corollary}
Every connected base with at least one edge and a unique maximum
independent set satisfies Theorem~\ref{thm:main}.
\end{corollary}
\begin{proof}
An outside vertex with one neighbour would give a different maximum
independent set by replacement. The zero-neighbour case contradicts
maximality.
\end{proof}
This graph-class inclusion is classical \cite[Lemma 2.3]{LMStructure};
the conclusion about six-spin obstructions is its application here.

The construction is effective for a finite input, without a
polynomial-time complexity claim. Enumerate maximum independent
sets and the leading kernel. Search through rational
$\delta=2^{-t}$, $t\ge3$, until exact congruence certifies
inertia $(2,2,0)$ for $U_\delta L U_\delta^\top$.
Then search through $\ell=-1+2^{-u}$, with $u>t+1$, until the exact
quadratic Hessian restriction has two positive eigenvalues.
The two limiting arguments prove termination of these unbounded
searches. Enumerating all occupations or subsets is generally
expensive; a practical search cap is not a mathematical bound.

For a direct check of the second step, sum over occupations $c$
of the base classes, with $\sum c_v=h$. Put
$b_v=s-c_v$,
$r_v=c_v+\sum_{u\sim v}c_u$,
$w(c)=\prod_v\binom{s}{c_v}$, and
\[
 e(c)=\sum_v\binom{c_v}{2}+\sum_{uv\in E(F)}c_uc_v.
\]
For witness spins $i,j$, let $u_i(v)=(1-\ell a_i)^{r_v}$ and
$v_i=\sum_vb_vu_i(v)$. The normalized Hessian entry is
\begin{align}
 \sum_c w(c)(1-\ell^2)^{e(c)}
 \Bigl[\,v_iv_j-\sum_vb_vu_i(v)u_j(v)
       -a_ia_j\,E_{ij}(c)\Bigr],                \label{eq:direct}\\
 E_{ij}(c)=
 \sum_vb_v(b_v-1)u_i(v)u_j(v)
 +\sum_{uv\in E(F)}b_ub_v
 \bigl(u_i(u)u_j(v)+u_i(v)u_j(u)\bigr).
 \nonumber
\end{align}
The omitted factor is $h!(|V(G)|-h-2)!>0$.
This formula counts outside ordered pairs directly and is suitable
for rational verification.

\section{Exact checks and reproducibility}

All 994 connected graphs of orders three through seven in the
NetworkX graph atlas were screened by two equivalent definitions:
edge deletion and the maximum-independent-set neighbour condition.
Exactly 360 satisfy the hypothesis. For each such base, both
$s=2$ and $s=3$ were checked, giving 720 blowups.
For every case, a separate enumeration of bounded occupation
vectors found all minimum-edge patterns. They agreed exactly with
Lemma~\ref{lem:ground}, including multiplicities, and the first
four-profile block agreed with Proposition~\ref{prop:block}.
These checks support the argument; the general theorem follows
from its proof, not an extrapolation of this census.

Four complete rational Hessian witnesses were also constructed
using \eqref{eq:direct}. All four used the same nodes
\[
 (a_0,\ldots,a_5)=
 \frac1{32}(-31,0,-28,-24,-20,-16).
\]
This is a statement about the four examples, not a universal
choice for the theorem.

\begin{center}
\begin{tabular}{@{}lrrrr@{}}
\toprule
Base $F$ & $s$ & $h$ & Occupations & Hessian inertia\\
\midrule
$P_3$ & 2 & 3 & 7 & $(2,2,0)$\\
$C_6$ & 2 & 5 & 126 & $(2,2,0)$\\
$K_4$ minus one edge & 2 & 3 & 16 & $(2,2,0)$\\
Mixed-degree example below & 2 & 5 & 51 & $(2,2,0)$\\
\bottomrule
\end{tabular}
\end{center}
The last base has vertex set $\{0,1,2,3,4\}$ and edges
$\{03,13,04,14,24,34\}$. Its unique maximum independent set is
$\{0,1,2\}$; outside degrees into that set are two and three.

The supplement includes the graph lists, occupation checks,
exact rational matrices and congruence records. The driver
\path{verify_all.py} reproduces both checks and compares all exact
records with the archived results. The witness-construction script
has finite practical search caps, documented separately
from the unbounded terminating construction above.
No floating-point eigenvalue calculation is used as a correctness
certificate. The mathematical checks are not a formal proof
assistant verification.

\section{Scope and remaining questions}

The result handles uniform clique blowups of $\alpha^-$-stable
bases. It does not assert that every noncomplete base is stable:
odd cycles, for example, have critical edges and fall outside the
hypothesis. It does not classify all graphs whose partition
polynomials are universally Lorentzian. The six-spin construction
also leaves open whether fewer spins suffice, whether target
parameters can be chosen uniformly for useful subclasses, and
whether effective parameter bounds can avoid full enumeration.

The proof uses a general transfer mechanism: a nonsingular initial
profile block can be realized by a small number of evaluation
nodes, with a separate base spin and pin spin. Applying that
mechanism requires an actual initial block with the needed
inertia; an arbitrary principal minor at later profiles does not
automatically give the same spin bound.

\paragraph{Assistance and review status.}
Substantial OpenAI Codex assistance was used for the argument,
writing and exact computations. The author operates AIRR.
The separate source argument received a technical audit by
GPT-6 Astra, requested at effort \texttt{xhigh}; that audit found no
material mathematical objection and included its own exact checks.
Version-specific manuscript assessments are recorded separately
with artifact hashes and must not be transferred between versions.
No independent human peer review, formal verification, or worldwide
priority determination is claimed. The numerical assessments of
earlier papers do not apply to this manuscript.

\begin{thebibliography}{9}
\bibitem{BH}
P. Br\"and\'en and J. Huh.
\emph{Lorentzian polynomials}.
\href{https://arxiv.org/abs/1902.03719}{arXiv:1902.03719}.

\bibitem{LOS}
J. Lee, J. Oh and J. Seo.
\emph{Counting homomorphisms in antiferromagnetic graphs via
Lorentzian polynomials}.
\href{https://arxiv.org/abs/2506.13659}{arXiv:2506.13659}, version 2.

\bibitem{LS}
J. Lee and J. Seo.
\emph{Antiferromagnetic models are clique-minimizing}.
\href{https://arxiv.org/abs/2608.17920}{arXiv:2608.17920}, version 1.

\bibitem{Eriksson}
L. Eriksson.
\emph{Rank-two obstructions to Lorentzian graph partition functions}.
AIRR, ARR-2026-7993ZVQWNC9E38P2, version 1, 17 October 2026.
\href{https://airr.science/papers/ARR-2026-7993ZVQWNC9E38P2/versions/v1/}
{Permanent version record}. Working paper.

\bibitem{LM}
V. E. Levit and E. Mandrescu.
\emph{Matrices and $\alpha$-Stable Bipartite Graphs}.
\href{https://arxiv.org/abs/math/9912120}{arXiv:math/9912120}, 1999.
Theorem 2.1 attributes the stability characterization to
T. W. Haynes, L. M. Lawson, R. C. Brigham and R. D. Dutton.

\bibitem{LMStructure}
V. E. Levit and E. Mandrescu.
\emph{On the structure of $\alpha$-stable graphs}.
\href{https://arxiv.org/abs/math/9911227}{arXiv:math/9911227}, 1999.

\bibitem{Turan}
P. Tur\'an.
\emph{On an extremal problem in graph theory} (in Hungarian).
Matematikai \'es Fizikai Lapok \textbf{48} (1941), 436--452.
\end{thebibliography}
\end{document}
